\documentclass{beamer}
\usepackage{amsmath}
\usepackage{amssymb}

% Choose a theme
\usetheme{Madrid}
\usecolortheme{default}
\usefonttheme{serif}

% --- CUSTOM FOOTER ---
% Redefine the footline to remove author/institute and keep title/page number
\makeatletter
\setbeamertemplate{footline}
{
  \leavevmode%
  \hbox{%
  \begin{beamercolorbox}[wd=.5\paperwidth,ht=2.25ex,dp=1ex,center]{title in head/foot}%
    \usebeamerfont{title in head/foot}\insertshorttitle
  \end{beamercolorbox}%
  \begin{beamercolorbox}[wd=.5\paperwidth,ht=2.25ex,dp=1ex,right]{date in head/foot}%
    \usebeamerfont{date in head/foot}\insertframenumber{} / \inserttotalframenumber\hspace*{2ex}
  \end{beamercolorbox}}%
  \vskip0pt%
}
\makeatother
% --- END CUSTOM FOOTER ---

\title[Lecture 14: Seismic sources]{Lecture 14: Seismic sources}
\author{David Al-Attar}
\institute{Michaelmas term 2024}
\date{}

\begin{document}

% --- INTRODUCTION ---

\begin{frame}
    \titlepage
\end{frame}

\begin{frame}{Outline and Motivation}
    We will now examine how seismic waves are generated by an earthquake.
    \begin{itemize}
        \item First, we'll see that our current equations of motion predict that earthquakes are impossible.
        \pause
        \item We will then modify the theory to account for displacement across a fault.
        \pause
        \item This is known as the \textbf{kinematic source approximation}, as the fault displacement is a specified parameter, not a dynamic variable.
        \pause
        \item Finally, this framework allows us to introduce the concept of \textbf{inverse problems}: using seismic recordings to determine what happened at the source.
    \end{itemize}
\end{frame}

% --- ELASTIC REBOUND ---

\begin{frame}{What Causes Earthquakes?}
    \begin{block}{The Source of Earthquakes}
        Most natural earthquakes are caused by motion across \textbf{faults} in the Earth. This motion is the result of a slow build-up and then rapid release of energy from long-term tectonic deformation.
    \end{block}
    \pause
    \begin{itemize}
        \item Most earthquakes occur near plate boundaries.
        \pause
        \item The idea that earthquakes are produced by rapid motion across a fault was developed by H.F. Reid after the 1909 San Francisco earthquake.
    \end{itemize}
\end{frame}

\begin{frame}{Reid's Elastic Rebound Theory}
    Reid's theory provides a simple conceptual model for the earthquake cycle.
    \pause
    \begin{enumerate}
        \item \textbf{Energy Storage:} Tectonic forces gradually build up elastic deformation on either side of a locked fault, storing potential energy in the rock.
        \pause
        \item \textbf{Rupture:} The stress across the fault eventually overcomes frictional resistance, causing a rapid rupture and slip.
        \pause
        \item \textbf{Energy Release:} Some of the stored elastic energy is radiated away as seismic waves.
    \end{enumerate}
    \pause
    \begin{alertblock}{The Kinematic Approximation}
        Instead of modeling the complex physics of rupture, we will treat the motion on the fault as a known parameter and calculate the waves it generates.
    \end{alertblock}
\end{frame}

% --- STRESS GLUT ---

\begin{frame}{A Problem With Our Theory}
    Our current theory for a hyperelastic body has a problem.
    \begin{itemize}
        \item The governing equations conserve the total energy of the body.
        $$ \frac{d}{dt}\int_{M} \left( \frac{1}{2}\rho||\mathbf{v}||^2 + W(\mathbf{F}) \right) d^3x = 0 $$
        \pause
        \item If the body is initially at rest in a stable equilibrium, any deformation would increase its potential energy.
        \pause
        \item Since kinetic energy cannot be negative, this requires an overall increase in total energy, which is forbidden by energy conservation.
    \end{itemize}
    \pause
    \begin{alertblock}{Conclusion}
        According to our theory of hyperelasticity, a body at rest must remain at rest forever. Earthquakes are impossible!
    \end{alertblock}
\end{frame}

\begin{frame}{Modifying the Theory for a Source}
    To allow for earthquakes, we must modify the potential energy.
    \begin{itemize}
        \item We must break the \textbf{time-translation invariance} of the system, because an earthquake has a specific onset time.
        \pause
        \item We do this by allowing the strain energy function $W$ to have an explicit time dependence: $W(\mathbf{x}, t, \mathbf{F})$.
        \pause
        \item We propose a specific form for this, separating the material's elastic properties ($U$) from the seismic source ($\overline{S}_{ij}$):
        $$ W(\mathbf{x}, t, \mathbf{F}) = U(\mathbf{x}, \mathbf{C}) - \frac{1}{2}\overline{S}_{ij}(\mathbf{x}, t)C_{ij} $$
    \end{itemize}
\end{frame}

\begin{frame}{The Stress Glut}
    The new term in the strain energy modifies the constitutive relation.
    \pause
    \begin{block}{Definition: The Stress Glut}
        The \textbf{stress glut} ($\overline{\sigma}_{ij}$) represents the difference between the stress predicted by the material's elastic properties and the true stress in the body.
        $$ \overline{\sigma}_{ij} = \underbrace{2F_{ik}\frac{\partial U}{\partial C_{kj}}}_{\text{Elastic Stress}} - \underbrace{T_{ij}}_{\text{True Stress}} $$
    \end{block}
    \pause
    \begin{exampleblock}{Physical Meaning}
        The stress glut can be interpreted as a failure of the elastic constitutive relation. This is consistent with the idea that material near a fault does not behave elastically during an earthquake.
    \end{exampleblock}
\end{frame}

\begin{frame}{Linearised Equations with a Source}
    Including the time-dependent source term in our linearised theory yields a new equation of motion.
    \begin{alertblock}{Linearised Equation with a Source}
        The seismic source enters the equation as an effective body force:
        $$ \rho\frac{\partial^2 u_{i}}{\partial t^2} - \frac{\partial}{\partial x_{j}}\left(A_{ijkl}\frac{\partial u_{k}}{\partial x_{l}}\right) = -\frac{\partial\overline{S}_{ij}}{\partial x_{j}} $$
    \end{alertblock}
    \pause
    \begin{block}{Boundary Condition with a Source}
        It also acts as an effective surface traction on the boundary:
        $$ A_{ijkl}\frac{\partial u_{k}}{\partial x_{l}}\hat{n}_{j} = \overline{S}_{ij}\hat{n}_{j} $$
    \end{block}
    \small
    Note: In the linearised theory, we use $\overline{S}_{ij}$ to denote the stress glut.
\end{frame}

\begin{frame}{Modeling an Idealised Fault}
    We now find the form of the stress glut for a fault.
    \begin{itemize}
        \item We model the fault as a mathematical surface, $\Sigma$, across which the displacement is discontinuous.
        \pause
        \item This is a good approximation as long as the seismic wavelength is much larger than the width of the actual fault zone.
        \pause
        \item We can write the displacement field $\mathbf{u}$ as a continuous part $\mathbf{u}^0$ plus a jump $\mathbf{w}$ across the fault plane (here, at $x_3=0$) using the Heaviside step function $H(x_3)$:
        $$ u_{i}(\mathbf{x}, t) = u_{i}^{0}(\mathbf{x}, t) + w_{i}(x_1, x_2, t)H(x_3) $$
    \end{itemize}
\end{frame}

\begin{frame}{Deriving the Stress Glut for a Fault}
    \begin{itemize}
        \item When we differentiate the discontinuous displacement $u_i$ to find the strain, the derivative of the Heaviside function produces a Dirac delta function $\delta(x_3)$.
        $$ \frac{\partial u_{i}}{\partial x_{j}} = \frac{\partial u_{i}^{0}}{\partial x_{j}} + \dots + w_{i}\delta_{j3}\delta(x_3) $$
        \pause
        \item This delta function in the strain leads to a \textbf{singular}, non-physical term in the elastic stress.
    \end{itemize}
    \pause
    \begin{alertblock}{The Stress Glut for a Fault}
        We define the stress glut to be precisely this non-physical term, so that it cancels out from the true stress. For a general fault surface $\Sigma$ with normal vector $\mathbf{n}$, the result is:
        $$ \overline{S}_{ij} = A_{ijkl}w_{k}n_{l}\delta_{\Sigma} $$
    \end{alertblock}
\end{frame}

% --- INVERSE PROBLEMS ---

\begin{frame}{Forward vs. Inverse Problems}
    \begin{block}{The Forward Problem}
        Given a model of the Earth ($\rho, \mathbf{A}$) and a model of the earthquake source ($\mathbf{w}$), calculate the expected seismic waves (synthetic seismograms).
        $$ \text{Model Parameters} \quad \rightarrow \quad \text{Predicted Data} $$
    \end{block}
    \pause
    \begin{block}{The Inverse Problem}
        Given recorded seismic waves from seismometers, work backwards to determine the parameters of the model (e.g., what was the slip $\mathbf{w}$ on the fault?).
        $$ \text{Observed Data} \quad \rightarrow \quad \text{Model Parameters} $$
    \end{block}
\end{frame}

\begin{frame}{Well-Posed vs. Ill-Posed Problems}
    There is a fundamental difference between forward and inverse problems.
    \begin{alertblock}{Well-Posed Problems}
        A problem is considered \textbf{well-posed} if it satisfies three conditions:
        \begin{enumerate}
            \item \textbf{Existence:} A solution exists for any valid input.
            \item \textbf{Uniqueness:} The solution is unique for a given input.
            \item \textbf{Continuity (Stability):} The solution depends continuously on the input (small changes in input lead to small changes in output).
        \end{enumerate}
        A problem that fails any of these conditions is \textbf{ill-posed}.
    \end{alertblock}
\end{frame}

\begin{frame}{Why Inverse Problems are Hard}
    \begin{itemize}
        \item Any plausible physical theory, like our forward problem of wave propagation, must be well-posed. If it weren't stable, we could never test it experimentally due to measurement errors.
        \pause
        \item In contrast, geophysical inverse problems are almost always \textbf{ill-posed}.
        \pause
        \item The main reason is non-uniqueness. We are trying to determine continuous functions (like slip on a fault) from a finite number of measurements. The problem is fundamentally \textbf{underdetermined}.
    \end{itemize}
\end{frame}

% --- SUMMARY ---

\begin{frame}{Summary}
    \begin{itemize}
        \item Most natural earthquakes are generated by rapid movement across faults. Reid's \textbf{elastic rebound theory} provides a conceptual model for this process.
        \pause
        \item A seismic source can be introduced into elasticity by modifying the strain energy, which gives rise to a \textbf{stress glut}.
        \pause
        \item The stress glut for an idealised fault is proportional to the displacement discontinuity across the fault surface.
        \pause
        \item A \textbf{forward problem} maps model parameters to data; an \textbf{inverse problem} maps observed data to model parameters.
        \pause
        \item Inverse problems are typically \textbf{ill-posed} (non-unique), while forward problems representing physical laws must be \textbf{well-posed}.
    \end{itemize}
\end{frame}

\end{document}
